% block: preface-heading page:7 type:section
\section*{\BT{Preface}{序言}}\phantomsection\label{sec:preface}
% block: preface-purpose page:7 type:ordinary-paragraph
\BT{In many fields of scientific computing and data science, we frequently face the problem of solving large and sparse linear systems of the form $A\bm x=\bm b$, which is one of the most time-consuming parts of all computations. From this fact, many researchers have devoted themselves to developing efficient numerical algorithms for solving the linear systems, and Krylov subspace methods are nowadays popular numerical algorithms and are known as one of the top ten algorithms of the twentieth century, others including fast Fourier transform and Quick Sort \BookCite{39}. Though the basic theory was established in the twentieth century, Krylov subspace methods have been developed by mathematicians, engineers, physicists, and many others.}{在科学计算和数据科学的许多领域, 我们经常面临求解形如 $A\bm x=\bm b$ 的大型稀疏线性方程组的问题, 这是整个计算过程中最耗时的部分之一. 因此, 许多研究者致力于开发求解线性方程组的高效数值算法. 如今, Krylov 子空间方法是广泛使用的数值算法, 被认为是 20 世纪十大算法之一, 其他算法包括快速 Fourier 变换和快速排序 \BookCite{39}. 尽管基本理论在 20 世纪就已建立, 数学家, 工程师, 物理学家和许多其他领域的研究者仍不断发展着 Krylov 子空间方法.}
% block: preface-books page:7 type:ordinary-paragraph
\BT{There are many excellent books on Krylov subspace methods, including those by:}{关于 Krylov 子空间方法, 已有许多优秀著作, 包括下列作者的著作:}
\begin{itemize}
\item Owe Axelsson, 1994 \BookCite{10},
\item Richard Barrett et al., 1994 \BookCite{18},
\item Wolfgang Hackbusch, 1994 \BookCite{92},
\item Are Magnus Bruaset, 1995 \BookCite{29},
\item Rüdiger Weiss, 1996 \BookCite{203},
\item Anne Greenbaum, 1997 \BookCite{81},
\item Yousef Saad, 2003 \BookCite{151},
\item Henk A. van der Vorst, 2003 \BookCite{196},
\item Jörg Liesen \BT{and}{和} Zdeněk Strakoš, 2012 \BookCite{122},
\item Gérard Meurant \BT{and}{和} Jurjen Duintjer Tebbens, 2020 \BookCite{129}.
\end{itemize}
% block: preface-features pages:7,8 type:ordinary-paragraph
\BT{In \BookCite{129}, detailed historical notes of Krylov subspace methods are described, which form a masterpiece (around 700 pages) of Krylov subspace methods for non-Hermitian linear systems. The features of this book are listed as follows:}{文献 \BookCite{129} 详细记述了 Krylov 子空间方法的历史, 是一部关于非 Hermitian 线性方程组的 Krylov 子空间方法的巨著 (约 700 页). 本书的特点如下:}
\begin{enumerate}[label=(\arabic*)]
\item \BT{Many applications of linear systems from computational science and data science;}{线性方程组在计算科学和数据科学中的许多应用;}
\item \BT{Krylov subspace methods for complex symmetric linear systems such as the COCG method and the COCR method;}{复对称线性方程组的 Krylov 子空间方法, 如 COCG 方法和 COCR 方法;}
\item \BT{Krylov subspace methods for non-Hermitian linear systems such as the BiCR method, the GPBiCG method, and the (block) IDR($s$) method;}{非 Hermitian 线性方程组的 Krylov 子空间方法, 如 BiCR 方法, GPBiCG 方法和 (块) IDR($s$) 方法;}
\item \BT{Krylov subspace methods for shifted linear systems such as the shifted IDR($s$) method;}{移位线性方程组的 Krylov 子空间方法, 如移位 IDR($s$) 方法;}
\item \BT{Matrix functions as applications of shifted linear systems.}{作为移位线性方程组应用的矩阵函数.}
\end{enumerate}
% block: preface-chapter-two page:8 type:ordinary-paragraph
\BT{Feature (1) corresponds to Chapter~\ref{chap:computational-applications}, and linear systems are derived from various applications: partial differential equations (finite difference discretization methods and the finite element method); computational physics: condensed matter physics (computation of Green's function) and lattice quantum chromodynamics (Wilson fermion matrix); machine learning (least-squares problems); matrix equations (Sylvester-type matrix equations); optimization (Hessian matrix over Euclidean space and Riemannian manifold using tensor computation notations).}{特点 (1) 对应于 Chapter~\ref{chap:computational-applications}, 其中从各种应用中导出线性方程组: 偏微分方程 (有限差分离散方法和有限元方法); 计算物理: 凝聚态物理 (Green 函数的计算) 和格点量子色动力学 (Wilson 费米子矩阵); 机器学习 (最小二乘问题); 矩阵方程 (Sylvester 型矩阵方程); 最优化 (使用张量计算记号表示 Euclidean 空间和 Riemannian 流形上的 Hessian 矩阵).}
% block: preface-chapter-three page:8 type:ordinary-paragraph
\BT{Features (2) and (3) correspond to Chapter~\ref{chap:krylov-classification}. In this chapter, Krylov subspace methods are classified into three groups: Hermitian linear systems, complex symmetric linear systems, and non-Hermitian linear systems. For Hermitian linear systems, the CG method is derived from the matrix form of the Lanczos process, the CR method is derived from the CG method, and the MINRES method is derived from the Lanczos process. For complex symmetric linear systems, the COCG method, the COCR method, and the QMR\_SYM method are described as extensions of the CG method, the CR method, and the MINRES method, respectively. For non-Hermitian linear systems, the BiCG method, the BiCR method, and the QMR method are described as extensions of the CG method, the CR method, and the MINRES method, respectively. The detailed derivations of the GPBiCG method and the (block) IDR($s$) method, one of the features of this book, are described in this chapter. In addition, some preconditioning techniques are briefly described.}{特点 (2) 和 (3) 对应于 Chapter~\ref{chap:krylov-classification}. 该章将 Krylov 子空间方法分成三类: 针对 Hermitian 线性方程组, 复对称线性方程组和非 Hermitian 线性方程组的方法. 对于 Hermitian 线性方程组, 从 Lanczos 过程的矩阵形式导出 CG 方法, 从 CG 方法导出 CR 方法, 并从 Lanczos 过程导出 MINRES 方法. 对于复对称线性方程组, 分别将 COCG 方法, COCR 方法和 QMR\_SYM 方法作为 CG 方法, CR 方法和 MINRES 方法的推广加以介绍. 对于非 Hermitian 线性方程组, 分别将 BiCG 方法, BiCR 方法和 QMR 方法作为 CG 方法, CR 方法和 MINRES 方法的推广加以介绍. 该章还详细推导了 GPBiCG 方法和 (块) IDR($s$) 方法, 这是本书的特点之一. 此外, 还简要介绍了一些预条件技术.}
% block: preface-chapter-four page:8 type:ordinary-paragraph
\BT{Feature (4) corresponds to Chapter~\ref{chap:shifted-linear-systems}. In this chapter, Krylov subspace methods for shifted linear systems are classified into three groups: Hermitian, complex symmetric, and non-Hermitian linear systems. The detailed derivations of these algorithms are described systematically.}{特点 (4) 对应于 Chapter~\ref{chap:shifted-linear-systems}. 该章将移位线性方程组的 Krylov 子空间方法分成三类: 针对 Hermitian, 复对称和非 Hermitian 线性方程组的方法, 并系统地给出这些算法的详细推导.}
% block: preface-chapter-five page:8 type:ordinary-paragraph
\BT{Feature (5) corresponds to Chapter~\ref{chap:matrix-functions}. If one needs a large matrix function, then Krylov subspace methods and Krylov subspace methods for shifted linear systems are methods of choice since these algorithms can produce any element of the matrix function. The definitions of matrix functions and well-known algorithms for matrix functions are also described.}{特点 (5) 对应于 Chapter~\ref{chap:matrix-functions}. 如果需要计算大型矩阵函数, 那么 Krylov 子空间方法以及移位线性方程组的 Krylov 子空间方法就是可选的方法, 因为这些算法能够计算矩阵函数中的任意元素. 该章还介绍矩阵函数的定义和一些著名的矩阵函数算法.}
% block: preface-numerical-experiments page:8 type:ordinary-paragraph
\BT{An additional feature of this book is that there are no numerical experiments except some typical numerical examples for further understanding the convergence behavior of Krylov subspace methods. The convergence of Krylov subspace methods depends highly on the coefficient matrix, and the best algorithm changes if the coefficient matrix changes. So, if the reader wants to solve linear systems, I recommend the reader to apply several Krylov subspace methods (including the BiCGSTAB method and the GMRES method) to their problem and choose the best one among them.}{本书的另一个特点是, 除了为进一步理解 Krylov 子空间方法的收敛行为而给出的一些典型数值例子以外, 没有数值实验. Krylov 子空间方法的收敛性高度依赖于系数矩阵, 系数矩阵改变时, 最佳算法也会随之改变. 因此, 如果读者需要求解线性方程组, 我建议对自己的问题尝试若干种 Krylov 子空间方法 (包括 BiCGSTAB 方法和 GMRES 方法), 并从中选出最佳方法.}
% block: preface-conclusion pages:8,9 type:ordinary-paragraph
\BT{This book is suitable for anyone who studied linear algebra and needs to solve large and sparse linear systems. I hope this book is helpful for the reader to understand the principles and properties of Krylov subspace methods and to correctly use Krylov subspace methods to solve their problems.}{本书适合任何学过线性代数并且需要求解大型稀疏线性方程组的读者. 我希望本书能够帮助读者理解 Krylov 子空间方法的原理和性质, 并正确地使用 Krylov 子空间方法求解自己的问题.}
\begin{flushright}
\BT{Nagoya, Japan}{日本名古屋}\qquad Tomohiro Sogabe\\
\BT{September 2022}{2022 年 9 月}
\end{flushright}
% block: acknowledgements-heading page:10 type:section
\section*{\BT{Acknowledgements}{致谢}}\phantomsection\label{sec:acknowledgements}
% block: acknowledgements-colleagues page:10 type:ordinary-paragraph
\BT{I would like to thank Prof. Dr. Lei Du at Dalian University of Technology for kindly sending me the LaTeX code of his block IDR($s$) method. I appreciate Masayuki Nakamura's work. I could raise the priority of writing this book from his timely reminders, together with his kind seasonal greetings. I thank Lokeshwaran Manickavasagam and Poojitha Ravichandran for their help with the production of this book. I am grateful to three anonymous reviewers for their careful reading of the manuscript and highly constructive comments that enhanced the quality of the manuscript. I would like to express my sincere thanks to three editors from the Springer Series in Computational Mathematics (SSCM) editorial board: Prof. Dr. Randolph E. Bank, Prof. Dr. Wolfgang Hackbusch, Prof. Dr. Josef Stoer, and Prof. Dr. Harry Yserentant for providing the author with fruitful comments from broad and distinguished professional viewpoints.}{感谢大连理工大学的 Prof. Dr. Lei Du 慷慨地向我提供他的块 IDR($s$) 方法的 LaTeX 代码. 感谢 Masayuki Nakamura 的工作. 他及时的提醒以及亲切的节日问候, 使我能够提高撰写本书的优先级. 感谢 Lokeshwaran Manickavasagam 和 Poojitha Ravichandran 在本书制作过程中提供的帮助. 感谢三位匿名审稿人仔细阅读书稿, 并提出极具建设性的意见, 提高了书稿质量. 我衷心感谢 Springer Series in Computational Mathematics (SSCM) 编委会的三位编辑: Prof. Dr. Randolph E. Bank, Prof. Dr. Wolfgang Hackbusch, Prof. Dr. Josef Stoer 和 Prof. Dr. Harry Yserentant, 他们以广博而杰出的专业视角为作者提供了富有成效的意见.}
% block: acknowledgements-family page:10 type:ordinary-paragraph
\BT{Finally, I appreciate Prof. Dr. Shao-Liang Zhang who introduced me to the field of Krylov subspace methods when I was a Ph.D. student at the University of Tokyo about 20 years ago, and I wish to thank my loving wife, Yuan Li, for her day-to-day support and her kindness. She recommended me to write this book on weekends, and she was taking care of two active sons: Kazuma and Takuma, and her old father.}{最后, 感谢 Prof. Dr. Shao-Liang Zhang, 大约 20 年前我在东京大学攻读博士学位时, 是他引领我进入 Krylov 子空间方法领域. 我还要感谢我亲爱的妻子 Yuan Li, 感谢她日复一日的支持与关爱. 她建议我利用周末撰写本书, 同时照料两个活泼的儿子 Kazuma 和 Takuma, 以及她年迈的父亲.}
\clearpage
